% Part: computability
% Chapter: computability-theory
% Section: equiv-ce-defs

\documentclass[../../../include/open-logic-section]{subfiles}

\begin{document}

\olfileid{cmp}{thy}{eqc}

\olsection[c.e. 集合的定义]{可计算枚举集的等价定义}


下面给出可计算枚举性的若干重要等价表述。

\begin{thm}
\ollabel{thm:ce-equiv}
设 $S$ 为自然数的集合。则以下各项等价：
\begin{enumerate}
\item\ollabel{case:ce} $S$ 是可计算枚举的。
\item\ollabel{case:ran-pc} $S$ 是某个\emph{部分}可计算函数的值域。
\item\ollabel{case:ran-prim} $S$ 是空集，或是某个原始递归函数的值域。
\item\ollabel{case:ce-domain} $S$ 是某个部分可计算函数的\emph{定义域}。
\end{enumerate}
\end{thm}

\begin{explain}
前三个条款表明，我们可以等价地将任何一个非空的可计算枚举集视为由某个可计算函数、部分可计算函数或原始递归函数所枚举的集合。第四个条款告诉我们，如果 $S$ 是可计算枚举的，那么对某个指标~$e$，
\[
S = \Setabs{x}{\cfind{e}(x) \fdefined}.
\]
换句话说，$S$ 就是使得计算 $\cfind{e}$ 会停机的那些输入所组成的集合。因此，可计算枚举集有时也被称为\emph{半可判定的}：如果一个数在集合中，你最终会得到“是”的回答；但如果不在，你永远等不到“否”的回答！{}
\end{explain}

\begin{proof}
由于每个原始递归函数都是可计算的，且每个可计算函数都是部分可计算的，因此 \olref{case:ran-prim} 蕴含 \olref{case:ce}，并且 \olref{case:ce} 蕴含~\olref{case:ran-pc}。（注意，如果 $S$ 是空集，那么 $S$~是那个处处无定义的部分可计算函数的值域。）如果我们能证明 \olref{case:ran-pc} 蕴含 \olref{case:ran-prim}，那么我们就证明了前三个条款彼此等价。

那么，假设 $S$ 是部分可计算函数 $\cfind{e}$ 的值域。如果 $S$ 是空集，证明已完成。否则，令 $a$ 为~$S$ 中的任意元素。根据 Kleene 范式定理，我们可以写出
\[
\cfind{e}(x) = U(\umin{s}{T(e, x, s)}).
\]
特别地，$\cfind{e}(x) \fdefined$ 且等于 $y$ 当且仅当存在某个 $s$ 使得 $T(e, x, s)$ 且 $U(s) = y$。定义 $f(z)$ 为
\[
f(z) = \begin{cases}
  U((z)_1) & \text{若 $T(e, (z)_0, (z)_1)$，} \\
  a        & \text{否则。}
\end{cases}
\]
那么 $f$ 是原始递归的，因为 $T$ 和 $U$ 是原始递归的。\iftag{TMs}{用Turing机的术语来说，如果 $z$ 编码了一个二元组 $\tuple{(z)_0, (z)_1}$，使得 $(z)_1$ 是机器~$M_e$ 在输入 $(z)_0$ 上的一个停机计算，则 $f$ 返回该计算的输出；否则，它返回~$a$。}{}

我们需要证明 $S$ 是~$f$ 的值域，即对任意自然数~$y$，$y \in S$ 当且仅当它在~$f$ 的值域中。正向推导：假设 $y \in S$。则 $y$ 在 $\cfind{e}$ 的值域中，因此存在某个 $x$ 和 $s$，使得 $T(e,x,s)$ 成立且 $U(s) = y$。那么 $y = f(\tuple{x,s})$。反之，假设 $y$ 在~$f$ 的值域中。则要么 $y = a$，要么存在某个~$z$，使得 $T(e,(z)_0,(z)_1)$ 且 $U((z)_1) = y$。由于在后一种情况下 $\cfind{e}(x) \fdefined = y$，因此无论哪种情形，$y$ 都在~$S$ 中。

（记号 $\cfind{e}(x) \fdefined = y$ 表示“$\cfind{e}(x)$ 有定义且等于 $y$”。我们本也可以直接写成 $\cfind{e}(x) = y$，但有时加上这个箭头有助于提醒我们处理的是偏函数。）

为了完成 \olref{thm:ce-equiv} 的证明，只需证明 \olref{case:ce} 和~\olref{case:ce-domain} 等价。首先，我们来证明 \olref{case:ce} 蕴含~\olref{case:ce-domain}。假设 $S$ 是某个可计算函数~$f$ 的值域，即
\[
S = \Setabs{y}{\text{存在某个 $x$，使得 } f(x) = y}.
\]
令
\[
g(y) = \umin{x}{(f(x) = y)}.
\]
那么 $g$ 是一个部分可计算函数，并且 $g(y)$ 有定义当且仅当存在某个~$x$，使得 $f(x) = y$。换言之，$g$ 的定义域就是~$f$ 的值域。\iftag{TMs}{用Turing机的术语来说：给定一个枚举 $S$ 中元素的Turing机~$F$，令 $G$ 为如下半判定 $S$ 的Turing机：它搜索~$F$ 的输出以检查给定元素是否在集合中，若在则停机，否则永远搜索下去。}{}

最后，为了证明 \olref{case:ce-domain} 蕴含~\olref{case:ce}，假设 $S$~是部分可计算函数~$\cfind{e}$ 的定义域，即
\[
S = \Setabs{x}{\cfind{e}(x) \fdefined}.
\]
如果 $S$ 是空集，证明已完成；否则，令 $a$ 为 $S$ 中的任意元素。定义 $f$ 为
\[
f(z) = \begin{cases}
(z)_0 & \text{若 $T(e,(z)_0,(z)_1)$，} \\
a & \text{否则。}
\end{cases}
\]
那么，与上文类似，一个数 $x$ 在~$f$ 的值域中当且仅当 $\cfind{e}(x) \fdefined$，即当且仅当 $x \in S$。\iftag{TMs}{用Turing机的术语来说：给定一台半判定 $S$ 的机器 $M_e$，通过遍历所有可能的Turing机计算来枚举 $S$ 的元素，并返回对应于停机计算的那些输入。}{}
\end{proof}

\olref{thm:ce-equiv} 的第 \olref{case:ce-domain} 款为我们提供了一种枚举可计算枚举集的便捷方法：对每个~$e$，令 $W_e$ 表示 $\cfind{e}$ 的定义域，即
\[
W_e = \Setabs{x}{\cfind{e}(x) \fdefined}.
\]
那么，若 $A$ 是任一可计算枚举集，则存在某个~$e$，使得 $A = W_e$。

下面给出可计算枚举集的另一种刻画。

\begin{thm}
\ollabel{thm:exists-char}
一个集合 $S$ 是可计算枚举的，当且仅当存在一个可计算关系 $R(x,y)$，使得
\[
S = \Setabs{ x }{ \lexists[y][R(x,y)] }.
\]
\end{thm}

\begin{proof}
正向推导：假设 $S$ 是可计算枚举的。则存在某个 $e$，使得 $S = W_e$。对这个 $e$，我们可以将 $S$ 写为
\[
S = \Setabs{ x }{ \lexists[y][T(e, x, y)] }.
\]
反向推导：假设 $S = \Setabs{ x }{
  \lexists[y][R(x, y)] }$。定义 $f$ 为
\[
f(x) \simeq \umin{y}{R(x, y)}.
\]
那么 $f$ 是部分可计算的，且 $S$ 是 $f$ 的定义域。
\end{proof}


\end{document}